% This is text associated with the open textbook "Open Electromagnetics"
% See immediately below for author, date
% This work is licensed under the Creative Commons Attribution-ShareAlike 4.0 International License. (CC BY-SA 4.0) 
% To view a copy of the license, visit http://creativecommons.org/licenses/by-sa/4.0/

%ORB TITLE Characteristic Impedance
%ORB AUTHOR 0        # S.W. Ellingson, Virginia Tech (ellingson@vt.edu)
%ORB DATE 20180802   # yyyymmdd
%ORB ABSTRACT Characteristic Impedance

\label{m0052_Characteristic_Impedance} % <-- Must match name of this file and module directory name.  Don't mess with this!
\index{characteristic impedance|(}

Characteristic impedance is the ratio of voltage to current for a wave that is propagating in single direction on a transmission line.
This is an important parameter in the analysis and design of circuits and systems using transmission lines.
In this section, we formally define this parameter and derive an expression for this parameter in terms of the equivalent circuit model introduced in Section~\ref{m0029_Lumped-Element_Model}. 

Consider a transmission line aligned along the $z$ axis.
Employing some results from Section~\ref{m0027_Wave_Equations_for_a_Transmission_Line}, recall that
the phasor form of the wave equation in this case is
%
\begin{equation}
\frac{\partial^2}{\partial z^2} \widetilde{V}(z) -\gamma^2~\widetilde{V}(z) =0
\label{m0052_eWaveEqnV}
\end{equation}
%
where
%
\begin{equation}
\gamma \triangleq \sqrt{\left( R' + j\omega L' \right)\left( G' + j\omega C' \right)} 
\label{m0052_egamma}
\end{equation}
%
Equation~\ref{m0052_eWaveEqnV} relates the potential phasor $\widetilde{V}(z)$ to the equivalent circuit parameters $R'$, $G'$, $C'$, and $L'$.
An equation of the same form relates the current phasor $\widetilde{I}(z)$ to the equivalent circuit parameters:
%
\begin{equation}
\frac{\partial^2}{\partial z^2} \widetilde{I}(z) -\gamma^2~\widetilde{I}(z) =0
\label{m0052_eWaveEqnI}
\end{equation}
%
Since both $\widetilde{V}(z)$ and $\widetilde{I}(z)$ satisfy the \emph{same} linear homogeneous differential equation, they may differ by no more than a multiplicative constant.
Since $\widetilde{V}(z)$ is potential and $\widetilde{I}(z)$ is current, that constant can be expressed in units of impedance.
Specifically, this is the \emph{characteristic impedance}, so-named because it depends only on the materials and cross-sectional geometry of the transmission line -- i.e., things which determine $\gamma$ -- and not length, excitation, termination, or position along the line. 

To derive the characteristic impedance, first recall that the general solutions to Equations~\ref{m0052_eWaveEqnV} and \ref{m0052_eWaveEqnI} are
%
\begin{equation}
\widetilde{V}(z) = V_0^+ e^{-\gamma z} + V_0^- e^{+\gamma z}
\label{m0052_eV}
\end{equation}
%
\begin{equation}
\widetilde{I}(z) = I_0^+ e^{-\gamma z} + I_0^- e^{+\gamma z}
\label{m0052_eI}
\end{equation}
%
where $V_0^+$, $V_0^-$, $I_0^+$, and $I_0^-$ are complex-valued constants whose values are determined by boundary conditions; i.e., constraints on $\widetilde{V}(z)$ and $\widetilde{I}(z)$ at some position(s) along the line.  
Also, we will make use of the telegrapher's equations (Section~\ref{m0079_Telegraphers_Equations}):
%
%
\begin{equation}
-\frac{\partial}{\partial z} \widetilde{V}(z) = \left[ R' + j\omega L' \right]~\widetilde{I}(z)
\label{m0052_eTelegraphersEquation1p}
\end{equation}
%
\begin{equation}
-\frac{\partial}{\partial z} \widetilde{I}(z) = \left[ G' + j\omega C' \right]~\widetilde{V}(z)
\label{m0052_eTelegraphersEquation2p}
\end{equation}
%
We begin by differentiating Equation~\ref{m0052_eV} with respect to $z$, which yields
%
\begin{equation}
\frac{\partial}{\partial z}\widetilde{V}(z) = -\gamma\left[V_0^+ e^{-\gamma z} - V_0^- e^{+\gamma z}\right]
\end{equation}
%
Now we use this this to eliminate $\partial\widetilde{V}(z)/\partial z$ in Equation~\ref{m0052_eTelegraphersEquation1p}, yielding
%
\begin{equation}
\gamma\left[V_0^+ e^{-\gamma z} - V_0^- e^{+\gamma z}\right] = \left[ R' + j\omega L' \right]~\widetilde{I}(z)
\end{equation}
%
Solving the above equation for $\widetilde{I}(z)$ yields:
%
\begin{equation}
\widetilde{I}(z) = \frac{\gamma}{R' + j\omega L'}\left[V_0^+ e^{-\gamma z} - V_0^- e^{+\gamma z}\right]
\end{equation}
%
Comparing this to Equation~\ref{m0052_eI}, we note
%
\begin{equation}
I_0^+ = \frac{\gamma}{R' + j\omega L'}~V_0^+
\end{equation}
%
\begin{equation}
I_0^- = \frac{-\gamma}{R' + j\omega L'}~V_0^-
\end{equation}
%
We now make the substitution
%
\begin{equation}
Z_0 = \frac{R' + j\omega L'}{\gamma} 
\label{m0052_eZ0}
\end{equation}
%
and observe
%
\begin{equation}
\boxed{
\frac{V_0^+}{I_0^+} = \frac{-V_0^-}{I_0^-} \triangleq Z_0
}
\end{equation}
%
As anticipated, we have found that coefficients in the equations for potentials and currents are related by an impedance, namely, $Z_0$.  

Characteristic impedance can be written entirely in terms of the equivalent circuit parameters by substituting Equation~\ref{m0052_egamma} into Equation~\ref{m0052_eZ0}, yielding: 
%
\begin{equation}
\boxed{
Z_0 = \sqrt{\frac{R' + j\omega L'}{G' + j\omega C'}}
}
\end{equation}  
%
%ORB HIGHLIGHT
\begin{mdframed}[backgroundcolor=yellow!20,nobreak=true] 
The characteristic impedance 
\index{characteristic impedance!definition}
$Z_0$ ($\Omega$) is the ratio of potential to current in a wave traveling in a single direction along the transmission line.
\end{mdframed}

Take care to note that $Z_0$ is \emph{not} the ratio of $\widetilde{V}(z)$ to $\widetilde{I}(z)$ in general; rather, $Z_0$ relates only the potential and current waves traveling in the \emph{same} direction.

Finally, note that transmission lines are normally designed to have a characteristic impedance that is completely real-valued -- that is, with no imaginary component.  This is because the imaginary component of an impedance represents energy \emph{storage} (think of capacitors and inductors), whereas the purpose of a transmission line is energy \emph{transfer}.

%%%%%%%%%%%%%%%%%%%%%%%%

{\bf Additional Reading:}
\begin{itemize}
\item ``\href{https://en.wikipedia.org/wiki/Characteristic_impedance}{Characteristic impedance}'' on Wikipedia.
\end{itemize}

\index{characteristic impedance|)}
